% Part: second-order-logic
% Chapter: syntax-and-semantics
% Section: size-of-domain

\documentclass[../../../include/open-logic-section]{subfiles}

\begin{document}

\olfileid[hi]{sol}{syn}{siz}

\olsection{अनंत और \usetoken{S}{enumerable}
  \usetoken{P}{domain} का वर्णन}

समुच्चय~$M$ (डेडेकिंड) अनंत है तभी और केवल तभी जब कोई ऐसा !!a{injective}
फलन~$f\colon M \to M$ हो जो !!{surjective} न हो, अर्थात् जिसके लिए
$\ran{f} \neq M$। प्रथम-कोटि तर्कशास्त्र में हम एक-स्थानी
!!{function}~$f$ पर विचार कर सकते हैं और कह सकते हैं कि
!!a{structure}~$\Struct{M}$ में उसे दिया गया फलन~$\Assign{f}{M}$,
!!{injective} है और $\ran{f} \neq \Domain{M}$:
\[
\lforall[x][\lforall[y][(\eq[f(x)][f(y)] \to \eq[x][y])]] \land
\lexists[y][\lforall[x][\eq/[y][f(x)]]].
\]
यदि $\Struct{M}$ इस !!{sentence} को संतुष्ट करता है, तो
$\Assign{f}{M}: \Domain{M} \to \Domain{M}$ !!{injective} है, अतः
$\Domain{M}$ अनंत होना चाहिए। यदि $\Domain{M}$ अनंत है और इसलिए ऐसा फलन
अस्तित्व में है, तो हम $\Assign{f}{M}$ को वही फलन मान सकते हैं और
$\Struct{M}$ इस !!{sentence} को संतुष्ट करेगा। किंतु इसके लिए हमारी भाषा में
वह अतार्किक चिह्न~$f$ होना चाहिए जिसका हम इस प्रयोजन के लिए उपयोग करते हैं।
द्वितीय-कोटि तर्कशास्त्र में हम सीधे कह सकते हैं कि ऐसा फलन \emph{अस्तित्व
में है}। अब इसके लिए~$f$ आवश्यक नहीं है, और हमें शुद्ध द्वितीय-कोटि
तर्कशास्त्र का यह !!{sentence} मिलता है:
\[
\fn{Inf} \ident \lexists[u][(\lforall[x][\lforall[y][(\eq[u(x)][u(y)]
      \lif \eq[x][y])]] \land \lexists[y][\lforall[x][\eq/[y][u(x)]]])].
\]
$\Sat{M}{\fn{Inf}}$ तभी और केवल तभी जब $\Domain{M}$ अनंत है। तब हम
$\fn{Fin} \ident \lnot \fn{Inf}$ परिभाषित कर सकते हैं;
$\Sat{M}{\fn{Fin}}$ तभी और केवल तभी जब $\Domain{M}$ परिमित है। शुद्ध
प्रथम-कोटि तर्कशास्त्र का कोई एक !!{sentence} यह व्यक्त नहीं कर सकता कि
!!{domain} अनंत है, यद्यपि ऐसे वाक्यों का कोई अनंत समुच्चय ऐसा कर सकता है।
शुद्ध प्रथम-कोटि तर्कशास्त्र के !!{sentence} का ऐसा कोई समुच्चय नहीं है जो
किसी !!a{structure} में तभी और केवल तभी संतुष्ट हो जब उसका प्रांत परिमित हो।

\begin{prop}
$\Sat{M}{\fn{Inf}}$ तभी और केवल तभी जब $\Domain{M}$ अनंत है।
\end{prop}

\begin{proof}
$\Sat{M}{\fn{Inf}}$ तभी और केवल तभी जब
  $\Sat{M}{\lforall[x][\lforall[y][(\eq[u(x)][u(y)] \lif \eq[x][y])]]
    \land \lexists[y][\lforall[x][\eq/[y][u(x)]]]}[s]$, किसी~$s$ के लिए।
  यदि ऐसा है, तो $s(u)$ !!a{injective} फलन है और कोई $y \in \Domain{M}$,
  ~$s(u)$ के परिसर में नहीं है। इसके विपरीत, यदि कोई !!a{injective}
  $f\colon \Domain{M} \to \Domain{M}$ ऐसा है कि $\ran{f} \neq
  \Domain{M}$, तो $s(u) = f$ ऐसा ही चर-नियतन है।
\end{proof}

समुच्चय $M$ !!{enumerable} है यदि उसके !!{element} का ऐसा परिगणन हो:
\[
m_0, m_1, m_2, \dots
\]
(बिना पुनरावृत्ति के, किंतु संभवतः परिमित)। ऐसा परिगणन तभी और केवल तभी
अस्तित्व में है जब कोई !!a{element} $z \in M$ और फलन $f\colon M \to M$ ऐसा
हो कि $z$, $f(z)$, $f(f(z))$, \dots, ही~$M$ के सभी !!{element} हों। यदि
परिगणन अस्तित्व में है, तो $z = m_0$ और $f(m_k) = m_{k+1}$ (अथवा यदि $m_k$
परिगणन का अंतिम !!{element} है तो $f(m_k) = m_k$) अपेक्षित !!{element} और
फलन हैं। दूसरी ओर, यदि ऐसे $z$ और $f$ अस्तित्व में हैं, तो $z$, $f(z)$,
$f(f(z))$, \dots, ~$M$ का परिगणन है और $M$ !!{enumerable} है। हम
द्वितीय-कोटि तर्कशास्त्र में $z$ और~$f$ का अस्तित्व व्यक्त करके ऐसा
!!{sentence} बना सकते हैं जो किसी !!a{structure} में तभी और केवल तभी सत्य
हो जब वह !!{structure}, !!{enumerable} हो:
\[
\fn{Count} \ident
\lexists[z][\lexists[u][\lforall[X][((X(z) \land
      \lforall[x][(X(x) \lif X(u(x)))]) \lif \lforall[x][X(x)])]]]
\]

\begin{prop}
$\Sat{M}{\fn{Count}}$ तभी और केवल तभी जब $\Domain{M}$ !!{enumerable} है।
\end{prop}

\begin{proof}
मान लें $\Domain{M}$ !!{enumerable} है और $m_0$, $m_1$, \dots, उसका
परिगणन है। पुनरावृत्तियाँ हटाकर हम सुनिश्चित कर सकते हैं कि कोई $m_k$ दो
बार न आए। $f(m_k) = m_{k+1}$ परिभाषित करें और $s(z) = m_0$ तथा $s(u) = f$
मानें। हम दिखाते हैं कि
\[
\Sat{M}{\lforall[X][((X(z) \land \lforall[x][(X(x) \lif X(u(x)))])
    \lif \lforall[x][X(x)])]}[s]
\]
मान लें $M \subseteq \Domain{M}$ स्वेच्छ है। आगे मान लें कि
$\Sat{M}{(X(z) \land \lforall[x][(X(x) \lif
X(u(x)))])}[\Subst{s}{M}{X}]$। तब $\Subst{s}{M}{X}(z) \in M$, और जब भी
$x \in M$ हो, तब $(\Subst{s}{M}{X}(u))(x) \in M$ भी। दूसरे शब्दों में,
चूँकि $\varAssign{\Subst{s}{M}{X}}{s}{X}$, इसलिए $m_0 \in M$; और यदि
$x \in M$ तो $f(x) \in M$। अतः $m_0 \in M$, $m_1 = f(m_0) \in M$,
$m_2 = f(f(m_0)) \in M$, आदि। इसलिए $M = \Domain{M}$ और इस कारण
$\Sat{M}{\lforall[x][X(x)]}[\Subst{s}{M}{X}]$। चूँकि $M \subseteq
\Domain{M}$ स्वेच्छ था, प्रमाण पूरा हुआ:
$\Sat{M}{\fn{Count}}$.

अब मान लें कि $\Sat{M}{\fn{Count}}$, अर्थात्
\[
\Sat{M}{\lforall[X][((X(z) \land \lforall[x][(X(x) \lif X(u(x)))])
    \lif \lforall[x][X(x)])]}[s]
\]
किसी~$s$ के लिए। $m = s(z)$ और $f = s(u)$ मानें और $M = \{m, f(m),
f(f(m)), \dots\}$ पर विचार करें। इस प्रकार परिभाषित $M$ स्पष्टतः
!!{enumerable} है। तब
\[
\Sat{M}{(X(z) \land \lforall[x][(X(x) \lif X(u(x)))])
    \lif \lforall[x][X(x)]}[\Subst{s}{M}{X}]
\]
परिकल्पना से। साथ ही $\Sat{M}{X(z)}[\Subst{s}{M}{X}]$, क्योंकि $M \ni m =
\Subst{s}{M}{X}(z)$; और $\Sat{M}{\lforall[x][(X(x) \lif
X(u(x)))]}[\Subst{s}{M}{X}]$ भी, क्योंकि जब भी $x \in M$ हो, तब
$f(x) \in M$ भी। अतः पूर्वपक्ष और सप्रतिबंध, दोनों संतुष्ट हैं, इसलिए
उत्तरपक्ष भी संतुष्ट होना चाहिए:
$\Sat{M}{\lforall[x][X(x)]}[\Subst{s}{M}{X}]$। किंतु इसका अर्थ है कि
$M = \Domain{M}$; और चूँकि परिभाषा से $M$ !!{enumerable} है, इसलिए
$\Domain{M}$ भी !!{enumerable} है।
\end{proof}

\begin{prob}
!!{sentence} $\fn{Inf} \land \fn{Count}$ उन्हीं और केवल उन्हीं प्रांतों में
सत्य है जो !!{denumerable} हैं। $\fn{Count}$ की परिभाषा को इस प्रकार बदलें
कि वह कोई भिन्न !!{sentence} बन जाए जो सीधे यह व्यक्त करे कि प्रांत
!!{denumerable} है, और सिद्ध करें कि वह ऐसा करता है।
\end{prob}

\end{document}

